\documentclass{amsart}
\usepackage{amsmath,amssymb}
\usepackage{tikz}
\usepackage{cleveref}

\theoremstyle{plain}
\newtheorem{theorem}{Theorem}[section]
\newtheorem{lemma}[theorem]{Lemma}
\newtheorem*{remark}{Remark}

\begin{document}

\title{A Sample Paper}
\author{Test Author}
\maketitle

\begin{abstract}
This is a short abstract. See Section~\ref{sec:main} for details.
\end{abstract}

\section{Main results}
\label{sec:main}

We begin with a foundational theorem, rooted in classical work \cite{GKP1994,Stanley2011}.

\begin{theorem}
\label{thm:main}
For every integer $n \geq 0$, we have
\[
  \sum_{k=0}^{n} k = \frac{n(n+1)}{2}.
\]
This identity is well known.
\end{theorem}

\begin{proof}
The proof proceeds by induction on $n$.

Base case: $n=0$ yields $0=0$.

Inductive step: assume the identity for $n$; then
\begin{align}
  \sum_{k=0}^{n+1} k
    &= \sum_{k=0}^{n} k + (n+1) \\
    &= \frac{n(n+1)}{2} + (n+1) \\
    &= \frac{(n+1)(n+2)}{2}.
\end{align}
Compare with \cref{thm:main} and \ref{fig:diagram}; see also \ref{thm:missing}.
\end{proof}

\begin{figure}
\label{fig:diagram}
\begin{tikzpicture}
  \draw (0,0) -- (1,1);
\end{tikzpicture}
\caption{A simple diagram.}
\end{figure}

% An end-of-document comment line.

\begin{remark}
No label here; uses $$x^2 + y^2 = z^2$$ as display.
\end{remark}

\end{document}
